\section{相关工作}

数据驱动的工业时序故障诊断已形成成熟路线，其通常以预先定义的诊断任务、可用标签或可控的故障场景为基础。这些前提使模型性能能够通过分类或识别指标进行比较，但也限制了方法直接迁移到缺乏统一故障真值的现场历史档案。

核工程数据方法的范围远不止中子信号分析。一类工作依托物理模型、状态重建和先进数值模拟服务于反应堆监测与安全研究；CORTEX、McSAFER 和 METIS 等项目体现了这一建模与验证路线 \citep{demaziere2022cortex}。另一类工作面向核电厂运行中的在线监测与仪表通道评估，已形成明确的工程和维护支持基础 \citep{iaea2008olm1,iaea2008olm2,hashemian2011online}。在信号层面，中子噪声研究为非侵入式分析反应堆运行和中子测量系统相关现象提供了方法基础；围绕中子探测器验证和核电厂故障诊断的数据驱动研究，也已在预设信号构造或故障场景下开展 \citep{pazsit2010noise,tambouratzis2020neutron,tambouratzis2020decomposition,elbordany2024fault}。

通用语言代理与工具调用研究表明，语言模型可以规划动作、调用外部工具，并将工具输出纳入后续推理，而不必直接承担全部专业计算 \citep{yao2023react,schick2023toolformer,patil2024gorilla}。专业诊断系统进一步展示了这种分工的具体形态：ChatCAD 将专用 CAD 网络的分类、分割和报告结果转换为文本表示后交由语言模型整合 \citep{wang2024chatcad}；AgentMD 则将临床计算器构建为可执行工具，并由语言代理选择和调用 \citep{jin2025agentmd}。这些研究共同说明，专业模型或确定性工具可以保留对底层测量和计算的职责，而语言模型承担跨工具的信息组织。

在核工程中，这种职责划分还需要满足来源可追溯、技术术语受控和工程师复核等要求。现有核领域语言模型研究已探索诊断可靠性、预测性维护中的技术语言生成、核电厂文档解释，以及面向安全关键决策支持的检索增强框架 \citep{reeves2024reliability,lin2025cws,mapes2025uprate,ndum2026radiant}。这些工作为核工程中语言模型的可靠使用提供了直接基础，但其主要输入仍是技术文档、维护语境或已有的文本问题，而非无标签的原始仪表时序。

检索增强生成通过外部资料为知识密集型生成提供依据 \citep{lewis2020rag}。除端到端回答质量外，近年来的 RAG 评价也开始区分检索模块与生成模块的失效模式；例如 RAGChecker 将评价细化到检索和生成的多个环节 \citep{ru2024ragchecker}。这一视角对于核工程尤其重要：预先指定的重要来源是否仍进入检索结果前列、段落能否定位以及陈述是否受到证据支持，不能被单一相似度分数或语言流畅性所替代。现有核工程技术文本研究已关注检索、事实性、来源追溯和人工验证 \citep{reeves2024reliability,lin2025cws,mapes2025uprate,ndum2026radiant}，为安全相关场景中的证据定位、引用和人工复核提供了重要基础。

总体而言，现有研究分别覆盖了工业时序诊断、专业工具调用和技术文档 RAG；对于无统一故障真值的核仪表历史档案，如何将可追溯的测量侧观察与受来源策略约束的证据侧解释衔接起来，仍构成本文关注的研究缺口。
